% ******************************* Appendix A ***********************************

\chapter{Project Management Plan}
\label{appendix:pm}

This appendix describes how the team organised the 15 weeks of the project:
scope, milestones, risks, quality assurance, responsibilities, communication
and configuration management. Because the project is a research study, the plan
is built around experiments: every conclusion in this thesis must be traceable
to a logged run with a fixed seed and configuration.

% ------------------------------------------------------------------------------
\section{Scope and Estimation}

The work was divided into eight work packages (Table~\ref{tab:pm-wbs}). Time is
counted in working weeks W01--W15 of the 15-week plan;
Section~\ref{sec:pm-process} explains how these weeks map to the calendar.

\begin{table}[htbp]
\caption{Work breakdown}
\label{tab:pm-wbs}
\centering
\small
\begin{tabular}{p{0.22\textwidth}p{0.42\textwidth}p{0.13\textwidth}p{0.13\textwidth}}
\toprule
\textbf{Work package} & \textbf{Main content} & \textbf{Weeks} & \textbf{Owner} \\
\midrule
1. Survey and orientation & Paper reviews per area, literature matrix, research gaps, research questions, hypotheses & W01--W03 & Whole team \\
2. Data & ViVQA and ViTextVQA from official sources, fixed splits, train/test leakage checks, OCR cache & W02, W05, W08, W14 & SV2 \\
3. Baseline & Zero-shot Qwen2.5-VL-3B-Instruct (B1), reproduction of a published TextVQA score (R0), frozen inference configuration & W04--W07 & SV1 \\
4. QLoRA adaptation & QLoRA training pipeline; models B2, B3, B3-full, C3 and the shared adapter SH; LoRA configuration ablations & W06--W11 & SV3 \\
5. Evaluation & EM, VQA-Acc, ANLS, F1; paired McNemar tests; seed stability; full evaluation table & W05--W11 & SV4 \\
6. Analysis & Error analysis, question-type breakdown, POPE-vi hallucination, token budget, effect of normalization and word segmentation on metrics & W07, W11--W12, W15 & SV4, SV3 \\
7. Prototype & Web demo, yes/no threshold calibration, parity check against the evaluated system, demo video & W13--W14 & SV4 \\
8. Reporting and defence & Weekly reports, review slides, reproducibility package, thesis & Weekly, W15 & Whole team, coordinated by the leader \\
\bottomrule
\end{tabular}
\end{table}

The initial estimate in W01 assumed access to a university GPU server. In W06
the team confirmed that no such server was available, and costs were
re-estimated in Google Colab compute units (CU). From W08, before every GPU job
the team estimated GPU hours and CU (about 1.54 CU per hour on an L4 and 1.2 CU
per hour on a T4) and compared the estimate with the actual reading in Colab.
One full QLoRA training run on ViTextVQA takes about 12 L4 hours (B2: 11.99 h;
B3-full: 12.75 h), so the number of full training runs was the main constraint
of the plan.

% ------------------------------------------------------------------------------
\section{Management Objectives}

\begin{itemize}
  \item \textbf{Reproducibility:} every run logs split, seed, hardware,
  hyperparameters, training time and inference time.
  \item \textbf{No result selection after the fact:} the decision rule of every
  main experiment is committed to git before the experiment is run.
  \item \textbf{Resource constraint:} QLoRA only; every experiment runs on a
  single Colab GPU (T4 or L4).
  \item \textbf{Milestones:} the main research is finished before Review~2 and
  the faculty council.
\end{itemize}

By the end of W14, the \path{experiments/} directory contains 141 runs, each
with a log. Six official models are listed in
\path{configs/official_models.json}, which is generated automatically from
the logs, and all four adapters stored on Google Drive match their recorded
SHA-256 checksums.

% ------------------------------------------------------------------------------
\section{Risks}

The W01 plan listed eight risks (model too large, insufficient GPU, faulty data,
baseline hard to reproduce, poor OCR, training not converging, demo hard to
deploy, slow training). Table~\ref{tab:pm-risks} records the risks that actually
occurred and how the team handled them.

\begin{table}[htbp]
\caption{Risks that occurred and how they were handled}
\label{tab:pm-risks}
\centering
\small
\begin{tabular}{p{0.24\textwidth}p{0.08\textwidth}p{0.09\textwidth}p{0.48\textwidth}}
\toprule
\textbf{Risk} & \textbf{Week} & \textbf{Level} & \textbf{Response} \\
\midrule
No university GPU server & W06 & High & All training moved to Colab; upgraded to Colab Pro on 7 October 2026 to use the L4 \\
fp16 on the T4 corrupts answers & W06 & High & Dtype probe on 73 questions: 23 corrupted in fp16; switched to bf16 on the L4 (0/73 corrupted) \\
Data from the wrong source & W05 & High & The local data were ViVQA-X, not ViVQA, and ViTextVQA came from an unofficial mirror. ViVQA was downloaded again from the original repository, access to the gated ViTextVQA repository was requested, and the answers in \path{test_gt.json} were used \\
A full training run takes about 12 hours & W08 & Medium & The main comparison was first run on an 8,000-question subset (B3 versus B2-sub); the full B3-full run followed in W10 once the subset result was clear \\
Library API change & W06 & Low & \texttt{transformers} 5.17.0 removed \path{warmup_ratio}; replaced by \path{warmup_steps}, and library versions are logged \\
GPU PaddleOCR conflicts with PyTorch in the prototype & W13 & Medium & OCR runs in a separate process (\path{src/prototype/ocr_worker.py}) \\
The shared adapter is biased towards ``yes'' & W12--W13 & Medium & A threshold on the first-token logit, chosen on a calibration set disjoint from POPE-vi \\
EXIF rotation flags corrupt the OCR cache & W14 & Medium & Pre-registered audit (164 of 165 rotated images mismatched); \path{ocr.py} fixed, 167 images re-read, B3-full re-evaluated (EM 53.35 $\rightarrow$ 53.63) \\
Review schedule changed & W10 & Low & Review~2 moved to 29 October 2026; one cumulative slide deck is updated instead of rebuilt \\
\bottomrule
\end{tabular}
\end{table}

% ------------------------------------------------------------------------------
\section{Project Process}
\label{sec:pm-process}

The team works in weekly cycles. Each working week goes through five steps:

\begin{enumerate}
  \item \textbf{Plan:} choose the questions of the week from the overall plan and
  the previous week's results, and record them in the week's README.
  \item \textbf{Pre-register:} for each main experiment, write the decision rule
  (threshold, statistical test, dataset) and commit it to git before running.
  \item \textbf{Run:} code is written and self-tested on a local machine, then
  pushed to GitHub; Colab notebooks clone the repository and run on a GPU; logs
  and results are stored on Google Drive.
  \item \textbf{Verify:} numbers are recomputed from the raw prediction files
  instead of being copied by hand, and the repository is cloned again to verify
  each commit.
  \item \textbf{Close:} update the README, the week's team report and the
  related chapter drafts.
\end{enumerate}

Working weeks W01--W15 are numbered by the content of the 15-week plan, not by
the calendar. The team worked ahead of the calendar: on 2 November 2026
(calendar week~9) the team had reached W15. Calendar weeks start on Monday
7 September 2026. Table~\ref{tab:pm-weeks} maps the two.

\begin{table}[htbp]
\caption{Working weeks and calendar dates (2026)}
\label{tab:pm-weeks}
\centering
\small
\begin{tabular}{p{0.12\textwidth}p{0.55\textwidth}p{0.22\textwidth}}
\toprule
\textbf{Week} & \textbf{Main content} & \textbf{Calendar dates} \\
\midrule
W01--W02 & Orientation, 15-week plan, literature matrix, data analysis & 07--16 Sep \\
W03 & Review~1 package & 17 Sep (Review~1: 24 Sep) \\
W04 & First zero-shot baseline; baseline switched to Qwen2.5-VL-3B & 18--21 Sep \\
W05 & Official data; zero-shot baseline on both datasets & 22--27 Sep \\
W06 & fp16 diagnosis; QLoRA pipeline & 28--29 Sep \\
W07 & Configuration frozen on the L4; official B1; B2; R0 reproduction & 30 Sep--12 Oct \\
W08 & OCR-enhanced prompting; B3 on a subset & 12--14 Oct \\
W09 & Seed stability; $2\times2$ table; ViVQA with 2 epochs & 15--17 Oct \\
W10 & LoRA configuration ablations; B3-full & 19--21 Oct \\
W11 & Full evaluation table; shared adapter & 21--22 Oct \\
W12 & POPE-vi hallucination; error analysis & 23--24 Oct \\
W13 & Web prototype & 26--31 Oct \\
W14 & Code completion, demo video, reproducibility package & 01--02 Nov \\
W15 & Review~2 feedback, defence slides, thesis & from 02 Nov \\
\bottomrule
\end{tabular}
\end{table}

The faculty milestones follow the calendar (Table~\ref{tab:pm-milestones}). The
thesis grade consists of 50\% for the process, for which the team submits seven
reports to the supervisors, and 50\% for the defence.

\begin{table}[htbp]
\caption{Faculty milestones}
\label{tab:pm-milestones}
\centering
\small
\begin{tabular}{p{0.20\textwidth}p{0.24\textwidth}p{0.32\textwidth}p{0.14\textwidth}}
\toprule
\textbf{Milestone} & \textbf{Date} & \textbf{Content} & \textbf{Status} \\
\midrule
Review~1 & 24 Sep 2026 (calendar week 3) & Problem, gaps, research questions, plan; two lecturers other than the supervisors review & Done, 10 comments \\
Review~2 & 29 Oct 2026 (calendar week 8) & Reproduced baseline, method, preliminary results & Done \\
Faculty council & 30 Nov--06 Dec 2026 (calendar week 13) & Three members, 60 minutes per team & Upcoming \\
University council & 14--20 Dec 2026 (calendar week 15) & Five members & Upcoming \\
\bottomrule
\end{tabular}
\end{table}

% ------------------------------------------------------------------------------
\section{Quality Management}

In this project, quality first means that the numbers can be trusted. The team
used five measures:

\begin{itemize}
  \item \textbf{Pre-registered rules:} the threshold and conclusion rule of every
  main experiment are committed before results exist, for example the rule for
  the shared adapter (commit \texttt{324caeb}, W11) and the hallucination rule
  for H4 (commit \texttt{9209170}, W12).
  \item \textbf{Paired tests:} every comparison of two models on the same test
  set uses McNemar's test per question for EM and a paired bootstrap 95\%
  confidence interval for ANLS; the main result is repeated with seeds 42, 43
  and 44.
  \item \textbf{Recomputation from raw data:} the final results table
  (\path{src/eval/full_eval_table.py}) reads the \path{predictions.json}
  files directly instead of copying numbers from weekly reports.
  \item \textbf{Self-tests:} the main modules have a \texttt{--selftest} mode that
  runs on CPU, on both Windows and Linux, to catch errors before GPU time is
  spent.
  \item \textbf{Checksums and run signatures:} adapters, the OCR cache, data files
  and slides are recorded with SHA-256 checksums; each run carries a
  configuration signature, so results of a different configuration are never
  appended to the same run directory.
\end{itemize}

% ------------------------------------------------------------------------------
\section{Training Plan}

The project needs knowledge in four different areas. In W01 each member read at
least three papers in their own area and answered ten fixed questions for each
paper, including the mandatory question ``what can we use for our project?''
(Table~\ref{tab:pm-papers}).

\begin{table}[htbp]
\caption{Foundation papers by area}
\label{tab:pm-papers}
\centering
\small
\begin{tabular}{p{0.32\textwidth}p{0.58\textwidth}}
\toprule
\textbf{Area} & \textbf{Papers} \\
\midrule
SV1: VLM / Baseline & BLIP-2, Florence-2, Qwen2.5-VL \\
SV2: Dataset / OCR & ViVQA, ViTextVQA, TextVQA (LoRRA) \\
SV3: PEFT / Training & LoRA, QLoRA, VL-Adapter, LLaVA-1.5 \\
SV4: Evaluation / Deployment & POPE, VQAv2, survey of vision-language models (TPAMI 2024) \\
\bottomrule
\end{tabular}
\end{table}

The remaining skills were learned during the project: QLoRA training with PEFT
and bitsandbytes, PaddleOCR and VietOCR, McNemar's test, and the workflow between
local machines, GitHub, Colab and Google Drive. A bilingual glossary
(\path{docs/Glossary.md}) was maintained throughout.

% ------------------------------------------------------------------------------
\section{Deliverables}

Table~\ref{tab:pm-deliverables} lists the deliverables required by the project
outline, including the reproducibility package (status as of 10 November 2026).

\begin{table}[htbp]
\caption{Deliverables}
\label{tab:pm-deliverables}
\centering
\small
\begin{tabular}{p{0.38\textwidth}p{0.36\textwidth}p{0.16\textwidth}}
\toprule
\textbf{Deliverable} & \textbf{Location} & \textbf{Status} \\
\midrule
Source code: 46 Python modules (data, training, evaluation, prototype) & \path{src/} on GitHub (PLHGNAOH/ViVQA-VLM) & Done \\
30 Colab notebooks & \path{notebooks/} & Done \\
Logs of 141 runs & \path{experiments/} & Done \\
QLoRA adapters of the official models & Google Drive; checksums in \path{configs/official_models.json} & Done \\
Reproducibility package: guide, library versions, inference script & \path{docs/REPRODUCE.md}, \path{requirements-colab.txt}, \path{src/prototype/infer.py} & Done \\
Web prototype and demo video (10 min 28 s) & \path{src/prototype/}, Google Drive & Done \\
Weekly reports W01--W15 & \path{docs/W01/} to \path{docs/W15/} & W15 in progress \\
Review~1 and Review~2 slides (English, Vietnamese) & \path{docs/W03/}, \path{docs/04_slides/} & Done; defence slides in progress \\
Thesis & Overleaf & In progress \\
\bottomrule
\end{tabular}
\end{table}

% ------------------------------------------------------------------------------
\section{Responsibilities}

The team has four students, each responsible for one technical area of the
project outline (Table~\ref{tab:pm-roles}). Pham Le Huy Hoang is the team
leader: in addition to PEFT and training, he coordinates progress and integrates
the results of all areas. Cross-area work, for example B3, which needs the OCR
cache of SV2 and the training pipeline of SV3, is owned by the member who
produces the main result and reviewed by the leader. Individual contributions
per week are recorded in the weekly team reports.

\begin{table}[htbp]
\caption{Roles and responsibilities}
\label{tab:pm-roles}
\centering
\small
\begin{tabular}{p{0.20\textwidth}p{0.18\textwidth}p{0.32\textwidth}p{0.20\textwidth}}
\toprule
\textbf{Member} & \textbf{Role} & \textbf{Main responsibilities} & \textbf{Main directories} \\
\midrule
Ho Dac Toan (SE183249) & SV1: VLM / Baseline & Qwen2.5-VL architecture, zero-shot baseline reproduction, comparison with published scores & \path{src/models/}, \path{src/eval/run_baseline.py} \\
Tran Quoc Huan (SE193642) & SV2: Dataset / OCR & Official data, splits, leakage checks, OCR cache with PaddleOCR and VietOCR & \path{data/}, \path{src/data/} \\
Pham Le Huy Hoang (SE182670) & SV3: PEFT / Training; team leader & QLoRA training, LoRA configuration ablations, shared adapter; coordination and integration & \path{src/adaptation/}, \path{configs/}, \path{adapters/} \\
Van Quang Truong (SE182458) & SV4: Evaluation / Deployment & Metrics, statistical tests, hallucination, prototype and demo & \path{src/eval/}, \path{src/prototype/} \\
\bottomrule
\end{tabular}
\end{table}

% ------------------------------------------------------------------------------
\section{Communication}

\begin{table}[htbp]
\caption{Communication channels}
\label{tab:pm-comm}
\centering
\small
\begin{tabular}{p{0.26\textwidth}p{0.48\textwidth}p{0.16\textwidth}}
\toprule
\textbf{Channel} & \textbf{Purpose} & \textbf{Frequency} \\
\midrule
Meetings with the supervisors & Confirm major decisions, for example the choice of Qwen2.5-VL-3B as baseline on 7 October 2026 & At milestones and when needed \\
Reviews with external lecturers & External feedback; after each review the team fills in the AI\_Capstone\_Review123 checklist, emails it to the supervisors and reviewers, and uploads it to the shared folder & Twice \\
Weekly reports (README and team report in \path{docs/Wxx/}) & Work done, results, decisions and postponed items & Every working week \\
GitHub & Source code, small logs, documents; every commit has a description & Continuous \\
Google Drive & Images, adapters, OCR cache, demo video, review recordings & Continuous \\
Department email & Review dates, announced only during the review week & At milestones \\
\bottomrule
\end{tabular}
\end{table}

% ------------------------------------------------------------------------------
\section{Configuration Management}

\paragraph{Documents.} Weekly documents live in \path{docs/W01} to
\path{docs/W15}; each folder has a README that indexes the week and a team
report. Cross-week documents live in numbered folders: \path{00_admin}
(conventions, checklists), \path{01_proposal} (project outline),
\path{02_references} (literature matrix, reading notes), \path{03_report}
(chapter drafts) and \path{04_slides} (review slides). Each document has a
single live copy; weekly folders link to it instead of copying it. Old documents
are moved to \path{_archive/} instead of being deleted.

\paragraph{Source code.} Source code is managed with git in the public
repository \path{github.com/PLHGNAOH/ViVQA-VLM}, working on the
\texttt{main} branch; an \texttt{ai-context} branch is generated by a GitHub
Action to collect reports. Data, images and adapter weights are not committed
(blocked by \path{.gitignore}) but stored on Google Drive; their locations and
checksums are recorded in the logs and in \path{configs/official_models.json}.
Run names follow \path{docs/00_admin/CONVENTIONS.md} and encode the week,
method, model, dataset, number format, GPU and seed, for example
\path{W07_B2_qlora_r16_qwen25vl_vitextvqa_ep1_bf16_L4_seed42}. The
base model is pinned to a fixed HuggingFace revision.

\begin{table}[htbp]
\caption{Tools and infrastructure}
\label{tab:pm-tools}
\centering
\small
\begin{tabular}{p{0.26\textwidth}p{0.64\textwidth}}
\toprule
\textbf{Category} & \textbf{Tools} \\
\midrule
Languages and libraries & Python, PyTorch, HuggingFace Transformers, PEFT, bitsandbytes, \path{qwen_vl_utils} \\
OCR and Vietnamese text & PaddleOCR, VietOCR; pyvi for word segmentation \\
Training and evaluation & Google Colab Pro: NVIDIA L4 for training and inference, Tesla T4 for OCR (free T4 in W04--W06) \\
Local machines & Windows 11, 4 GB GPU: development, self-tests, preprocessing, error analysis \\
Storage & Google Drive (data, adapters, large logs); GitHub (code, small logs, documents) \\
Prototype & Gradio \\
Code and writing & Cursor IDE; AI assistants (Claude, Gemini) for code and drafting, with every number checked against the logs \\
Slides and thesis & Canva, PowerPoint, Google Docs, Overleaf \\
\bottomrule
\end{tabular}
\end{table}

The Colab budget was tracked in CU and recorded in the weekly reports. Remaining
CU read in Colab: 146.75 (13 Oct), 86.52 (19 Oct), 49.3 (21 Oct), 34.96
(23 Oct) and 17.38 (31 Oct). W10 alone, with the LoRA configuration ablations
and B3-full, used about 37 CU.
